\chapter{Antecedentes y estado del arte}
\label{chap:antecedentes}

\drop{L}{a calidad} de datos se ha consolidado como un campo de estudio fundamental
dentro de la informática y la gestión de la información. A lo largo de las
últimas tres décadas, la comunidad científica ha producido marcos conceptuales,
estándares internacionales y herramientas especializadas que buscan
sistematizar la evaluación y la mejora de los datos utilizados en sistemas de
información. En el contexto de la \ac{IA} y el \ac{ML}, la calidad de los datos
de entrenamiento condiciona directamente la fiabilidad de los modelos
resultantes, lo que ha impulsado una línea de investigación conocida como
\emph{data-centric AI}. Este capítulo examina los fundamentos teóricos, el marco
normativo, las herramientas existentes y la justificación académica que
sustentan el desarrollo de \dataqual como plataforma de evaluación de calidad de
datos para proyectos de \ac{IA}.


%% =========================================================================
\section{La importancia de la calidad de datos en proyectos de IA}
\label{sec:importancia-calidad}
%% =========================================================================

\subsection{Marcos conceptuales fundacionales}

Los trabajos seminales de Wang y Strong~\cite{wang1996} establecieron un marco
conceptual para la calidad de datos basado en cuatro categorías fundamentales:

\begin{itemize}
  \item \textbf{Calidad intrínseca:} precisión, objetividad, credibilidad y
    reputación de los datos.
  \item \textbf{Calidad contextual:} relevancia, valor añadido, completitud,
    cantidad apropiada y oportunidad temporal de los datos respecto a la tarea
    en la que se utilizan.
  \item \textbf{Calidad representacional:} interpretabilidad, facilidad de
    comprensión, representación concisa y representación consistente.
  \item \textbf{Accesibilidad:} disponibilidad de los datos y seguridad de
    acceso a los mismos.
\end{itemize}

Este marco ha sido ampliamente adoptado tanto en la investigación como en la
industria, y constituye la base sobre la que se han construido estándares
internacionales como \ac{ISO}/IEC 25012.

\subsection{Impacto económico de la mala calidad de datos}

Redman~\cite{redman1998, redman2001} cuantificó por primera vez el impacto
económico de la mala calidad de datos en las organizaciones, estableciendo que
las empresas gastan entre el 10\% y el 25\% de sus ingresos en gestionar
problemas derivados de datos de baja calidad. Estos hallazgos han sido
corroborados y ampliados por estudios posteriores del MIT Information Quality
Program y del \ac{DMBOK}. Según estimaciones de Gartner~\cite{gartner2017}, la
mala calidad de datos cuesta a las organizaciones una media de 12,9 millones de
dólares anuales.

En el contexto específico de proyectos de \ac{IA}, este coste se amplifica de
forma considerable: un modelo entrenado con datos defectuosos no solo produce
resultados incorrectos, sino que las decisiones basadas en esas predicciones
pueden tener consecuencias financieras, legales y reputacionales significativas.
Sambasivan et al.~\cite{sambasivan2021} demostraron, en un estudio realizado en
Google Research, que los problemas de calidad de datos constituyen la causa
principal de fallos en sistemas de \ac{IA} en producción, superando incluso a
los errores algorítmicos. Su trabajo, titulado «Everyone wants to do the model
work, not the data work», evidenció la brecha entre la atención que reciben los
modelos y la que reciben los datos que los alimentan.

Múltiples encuestas en la industria~\cite{kaggle2022} coinciden en que los
profesionales de datos dedican entre el 60\% y el 80\% de su tiempo a tareas
de limpieza y preparación de datos, lo que refleja tanto la prevalencia de los
problemas de calidad como la ausencia de herramientas automatizadas que
agilicen el diagnóstico.

\subsection{El paradigma data-centric AI}

Andrew Ng~\cite{ng2021datacentric}, uno de los investigadores más influyentes
en \ac{IA}, ha popularizado el concepto de \emph{data-centric AI} como
contraposición al enfoque tradicional \emph{model-centric}. Mientras que el
enfoque \emph{model-centric} se centra en mejorar los algoritmos manteniendo
los datos fijos, el enfoque \emph{data-centric} propone mantener los modelos
fijos y mejorar sistemáticamente los datos. Este cambio de paradigma reconoce
que, para la mayoría de aplicaciones prácticas de \ac{IA}, la mejora de los
datos produce un impacto mayor que la mejora de los algoritmos.

Diversas iniciativas han consolidado este paradigma en la comunidad:

\begin{itemize}
  \item La \emph{Data-Centric AI Competition}, organizada por Andrew Ng y
    Landing AI, en la que los participantes compiten mejorando un dataset fijo
    en lugar de un modelo fijo.
  \item \emph{DataPerf}, una iniciativa de MLCommons que proporciona
    \emph{benchmarks} centrados en la calidad de datos para sistemas de
    \ac{ML}.
  \item El área de investigación \emph{Data Quality for AI} (DQAI), que
    explora métricas de calidad específicas para conjuntos de datos destinados
    al entrenamiento de modelos.
\end{itemize}

Además, la problemática del sesgo algorítmico refuerza la necesidad de
herramientas de evaluación de calidad. Los datos desequilibrados o no
representativos son la raíz de muchos casos documentados de sesgo en sistemas de
\ac{IA}: desde algoritmos de contratación que discriminan por
género~\cite{amazon2018} hasta sistemas de reconocimiento facial con tasas de
error desproporcionadas para minorías étnicas~\cite{buolamwini2018}. La calidad
de los datos \textemdash en particular, el equilibrio de clases y la
representatividad\textemdash~es un factor determinante en la equidad de los
sistemas inteligentes.


%% =========================================================================
\section{Marco normativo: estándares ISO}
\label{sec:marco-normativo}
%% =========================================================================

La evaluación de la calidad de datos se enmarca dentro de un conjunto de
estándares internacionales publicados por la \ac{ISO}. \dataqual ha sido
diseñado para alinearse con tres estándares clave: \ac{ISO}/IEC~25012,
\ac{ISO}/IEC~25024 e \ac{ISO}/IEC~5259.

\subsection{ISO/IEC 25012: Modelo de calidad de datos}
\label{sec:iso25012}

El estándar \ac{ISO}/IEC 25012~\cite{iso25012} forma parte de la familia
\ac{SQuaRE} y define un modelo de calidad de datos que comprende 15
características organizadas en dos categorías principales:

\begin{itemize}
  \item \textbf{Calidad inherente:} propiedades que pertenecen a los datos
    en sí mismos, independientemente del sistema que los almacena. Incluye
    exactitud, completitud, consistencia, credibilidad y actualidad
    (\emph{currentness}).
  \item \textbf{Calidad dependiente del sistema:} propiedades que dependen
    del contexto tecnológico en el que se utilizan los datos. Incluye
    disponibilidad, portabilidad y recuperabilidad.
\end{itemize}

\dataqual implementa métricas que cubren las principales características
inherentes del estándar, añadiendo además métricas específicas para proyectos
de \ac{IA} y \ac{ML} que no están contempladas en la norma. El
Cuadro~\ref{tab:iso25012-dataqual} muestra la correspondencia entre las
características de \ac{ISO}/IEC 25012 y las métricas implementadas en
\dataqual.

\begin{table}[htbp]
  \centering
  \caption{Correspondencia entre características de ISO/IEC 25012 y métricas de \dataqual}
  \label{tab:iso25012-dataqual}
  \begin{tabularx}{\textwidth}{l l X}
    \toprule
    \textbf{Característica ISO 25012} & \textbf{Métrica \dataqual} & \textbf{Implementación} \\
    \midrule
    Completitud              & \texttt{completeness}        & Ratio de valores no nulos por columna \\
    Exactitud                & \texttt{syntactic\_accuracy}  & Conformidad con patrones de formato \\
    Consistencia (Con-ML-4)  & \texttt{logical\_consistency} & Reglas IF-THEN entre columnas \\
    Consistencia (Con-ML-1)  & \texttt{uniqueness}           & Detección de registros duplicados \\
    Actualidad               & \texttt{currentness}          & Antigüedad del dato más reciente \\
    \midrule
    \multicolumn{3}{l}{\emph{Métricas adicionales específicas para IA/ML}} \\
    \midrule
    ---             & \texttt{outliers}             & Detección de valores atípicos (\ac{IQR}, Z-score) \\
    ---             & \texttt{class\_balance}       & Equilibrio de clases (entropía de Shannon) \\
    \bottomrule
  \end{tabularx}
\end{table}

\subsection{ISO/IEC 5259: Calidad de datos para analítica e IA}
\label{sec:iso5259}

El estándar \ac{ISO}/IEC 5259~\cite{iso5259} es de reciente publicación y
aborda específicamente la calidad de datos en contextos de analítica e \ac{IA}.
Se estructura en cuatro partes:

\begin{enumerate}
  \item \textbf{Parte 1 \textemdash~Visión general:} establece los conceptos
    fundamentales y el vocabulario común.
  \item \textbf{Parte 2 \textemdash~Métricas de calidad de datos:} define
    métricas cuantitativas estandarizadas para evaluar la calidad de datos
    destinados a sistemas de \ac{IA}.
  \item \textbf{Parte 3 \textemdash~Gestión de datos:} cubre la gestión de
    conjuntos de datos, incluyendo versionado y trazabilidad.
  \item \textbf{Parte 4 \textemdash~Marco de procesos:} describe el flujo de
    evaluación integrado en el ciclo de vida de un proyecto de \ac{IA}.
\end{enumerate}

\dataqual se alinea con este estándar en varios aspectos: proporciona métricas
cuantitativas estandarizadas (Parte~2), implementa gestión de datasets con
versionado (Parte~3) y ofrece un flujo de evaluación integrado en el ciclo de
vida del proyecto (Parte~4).

Desde el punto de vista del alcance, la contribución más relevante de
\ac{ISO}/IEC 5259-2 es la organización de las características de calidad en
cuatro grupos según su naturaleza y su dependencia del contexto tecnológico:
\emph{características inherentes}, \emph{características mixtas} (inherentes y
dependientes del sistema), \emph{características dependientes del sistema} y
\emph{características adicionales} específicas para \ac{IA}/\ac{ML}. Esta
taxonomía, analizada en detalle en la Sección~\ref{sec:alcance-metodologico},
constituye el marco conceptual sobre el que se delimita el alcance técnico de
\dataqual.

\subsection{ISO/IEC 25024: Medición de la calidad de datos}
\label{sec:iso25024}

El estándar \ac{ISO}/IEC 25024~\cite{iso25024} complementa al 25012
proporcionando métricas concretas y fórmulas de medición para cada
característica de calidad. De este modo, se pasa de un modelo conceptual
abstracto a un conjunto de indicadores cuantificables.

\dataqual implementa directamente varias de las métricas definidas en este
estándar:

\begin{itemize}
  \item \textbf{Ratio de completitud:} proporción de valores no nulos sobre
    el total de valores esperados en una columna o dataset.
  \item \textbf{Ratio de unicidad:} proporción de valores únicos respecto al
    total de registros, utilizada para detectar duplicados.
  \item \textbf{Ratio de conformidad sintáctica:} proporción de valores que
    se ajustan a un patrón o formato predefinido (expresiones regulares,
    rangos numéricos, conjuntos de valores permitidos).
\end{itemize}

La adopción de estas fórmulas estandarizadas garantiza que los resultados de
\dataqual sean comparables con los obtenidos por otras herramientas que
implementen el mismo estándar, facilitando la interoperabilidad y la
auditoría.

Además, el Reglamento Europeo de Inteligencia
Artificial~\cite{aiact2024} exige que las organizaciones puedan demostrar la
calidad y trazabilidad de los datos utilizados en sus modelos, lo que convierte
la evaluación estandarizada de calidad de datos en una necesidad no solo técnica,
sino también de cumplimiento normativo.


%% =========================================================================
\section{Alcance metodológico: evaluación algorítmica de la calidad inherente}
\label{sec:alcance-metodologico}
%% =========================================================================

El marco normativo descrito en la sección anterior permite delimitar con
precisión qué dimensiones de calidad de datos son evaluables mediante el
análisis estático de un fichero de datos y cuáles dependen de factores externos
al propio dato. Esta distinción no es una limitación arbitraria de \dataqual,
sino una consecuencia directa de la naturaleza de las características definidas
por \ac{ISO}/IEC 5259-2~\cite{iso5259}. La presente sección justifica las
decisiones de alcance adoptadas en el proyecto a partir del marco conceptual
de dicho estándar.

\subsection{Taxonomía de características de calidad según ISO/IEC 5259-2}
\label{sec:taxonomia-iso5259}

\ac{ISO}/IEC 5259-2 organiza las características de calidad de datos para
analítica e \ac{IA} en cuatro grupos según su naturaleza y su grado de
dependencia del contexto tecnológico:

\begin{itemize}
  \item \textbf{Características inherentes} (5): exactitud
    (\emph{accuracy}), completitud (\emph{completeness}), consistencia
    (\emph{consistency}), credibilidad (\emph{credibility}) y actualidad
    (\emph{currentness}). Son propiedades intrínsecas del dato, evaluables
    a partir del análisis de sus valores sin necesidad de conocer el sistema
    que lo aloja.

  \item \textbf{Características mixtas} (7): accesibilidad
    (\emph{accessibility}), conformidad (\emph{compliance}), eficiencia
    (\emph{efficiency}), precisión (\emph{precision}), trazabilidad
    (\emph{traceability}), confidencialidad (\emph{confidentiality}) y
    comprensibilidad (\emph{understandability}). Dependen tanto de las
    propiedades del dato como del contexto del sistema que lo gestiona.

  \item \textbf{Características dependientes del sistema} (3):
    disponibilidad (\emph{availability}), portabilidad (\emph{portability})
    y recuperabilidad (\emph{recoverability}). Vienen determinadas
    exclusivamente por la infraestructura tecnológica de la organización.

  \item \textbf{Características adicionales para IA/ML} (9): equilibrio
    (\emph{balance}), diversidad (\emph{diversity}), efectividad
    (\emph{effectiveness}), identificabilidad (\emph{identifiability}),
    relevancia (\emph{relevance}), representatividad
    (\emph{representativeness}), similitud (\emph{similarity}), puntualidad
    (\emph{timeliness}) y auditabilidad (\emph{auditability}). Abordan
    necesidades específicas del aprendizaje automático que los marcos
    generalistas de calidad de datos no contemplan.
\end{itemize}

El Cuadro~\ref{tab:iso5259-cobertura} sintetiza la cobertura de \dataqual
respecto a las cuatro categorías anteriores.

\begin{table}[htbp]
  \centering
  \caption{Cobertura de \dataqual respecto al modelo de calidad de \ac{ISO}/IEC 5259-2}
  \label{tab:iso5259-cobertura}
  \small
  \begin{tabularx}{\textwidth}{l X l X}
    \toprule
    \textbf{Categoría} & \textbf{Características} & \textbf{\dataqual} & \textbf{Motivo} \\
    \midrule
    Inherentes
      & Exactitud, Completitud, Consistencia, Credibilidad, Actualidad
      & \checkmark~Implementadas
      & Extraíbles del dataset sin contexto externo \\[4pt]
    Mixtas
      & Accesibilidad, Conformidad, Eficiencia, Precisión, Trazabilidad,
        Confidencialidad, Comprensibilidad
      & Parcial
      & Requieren información del sistema que aloja el dato \\[4pt]
    Dep.\ del sistema
      & Disponibilidad, Portabilidad, Recuperabilidad
      & ---~Fuera de alcance
      & Determinadas por la infraestructura TI de la organización \\[4pt]
    Adicionales (MVP)
      & Equilibrio (\texttt{class\_balance}), Identificabilidad (enmascaramiento \ac{PII})
      & \checkmark~Implementadas
      & Impacto directo en fiabilidad y cumplimiento normativo de modelos de \ac{IA} \\[4pt]
    Adicionales (resto)
      & Diversidad, Efectividad, Relevancia, Representatividad, Similitud,
        Puntualidad, Auditabilidad
      & ---~Fuera de alcance
      & Requieren metadatos organizativos o definición de la población objetivo \\
    \bottomrule
  \end{tabularx}
\end{table}

\subsection{Características evaluables algorítmicamente}
\label{sec:caracteristicas-evaluables}

Las características inherentes de \ac{ISO}/IEC 5259-2 son evaluables
algorítmicamente porque sus métricas se calculan a partir de los propios
valores del dataset, sin necesidad de acceder a sistemas externos ni de conocer
el dominio de negocio. La completitud se mide como la proporción de valores no
nulos; la exactitud sintáctica, como la conformidad con patrones de formato
definidos; la consistencia, como la ausencia de contradicciones lógicas entre columnas
(sub-dimensión Con-ML-4 de \ac{ISO}/IEC 5259-2) y como la ausencia de registros
duplicados (sub-dimensión Con-ML-1); y la actualidad, como la antigüedad del valor
más reciente en columnas temporales.

Esta propiedad de \emph{auto-contención} es lo que hace posible desarrollar
una herramienta de auditoría que opera exclusivamente sobre el fichero de datos
sin requerir integración con los sistemas de la organización. El motor de
métricas de \dataqual explota esta característica al procesar cada
\texttt{AnalysisRun} de forma asíncrona y aislada, leyendo el dataset
almacenado en MinIO y calculando las métricas configuradas sin ninguna
conexión a los sistemas transaccionales del cliente.

\subsection{Características fuera del alcance técnico}
\label{sec:caracteristicas-fuera-alcance}

Dos categorías de características quedan fuera del alcance técnico de
\dataqual por razones de naturaleza estructural, no de elección de
implementación.

\paragraph{Dependencia de infraestructura organizativa.}
Las características dependientes del sistema ---disponibilidad, portabilidad
y recuperabilidad--- no describen propiedades del dato en sí, sino de la
infraestructura que lo almacena y sirve. La \emph{disponibilidad} mide si el
dataset puede ser recuperado por usuarios autorizados dentro de un plazo
acordado, lo que remite a acuerdos de nivel de servicio (\ac{SLA}) y
arquitecturas de almacenamiento. La \emph{portabilidad} evalúa la capacidad
de mover el dato entre sistemas preservando su calidad, lo que requiere acceso
a los conectores y formatos de exportación propietarios de la organización.
La \emph{recuperabilidad} depende de las políticas de copia de seguridad y los mecanismos
de restauración de datos. Ninguna de estas métricas puede calcularse analizando
el contenido de un fichero: exigen acceso al entorno operativo de la
organización, del que una herramienta de análisis de datos externo no dispone
por definición.

\paragraph{Dependencia de metadatos organizativos y contexto de negocio.}
Un subconjunto de las características adicionales para \ac{IA}/\ac{ML} requiere
información que trasciende el fichero de datos. La \emph{representatividad}
necesita la definición de la población objetivo del modelo, que es un
requisito de negocio no inscrito en el dato. La \emph{relevancia} y la
\emph{efectividad} requieren un diccionario de datos y políticas de gobernanza
privadas para determinar si el dataset es adecuado para la tarea prevista.
La \emph{diversidad} y la \emph{similitud} presuponen un conocimiento del
espacio de características del dominio que, desde una posición de análisis
externo, no resulta accesible sin la colaboración directa de la organización propietaria.

\subsection{Priorización de características para el MVP}
\label{sec:priorizacion-mvp}

Entre las características algorítmicamente evaluables, el desarrollo del
\ac{MVP} de \dataqual ha priorizado aquellas con mayor impacto en la fiabilidad
de los modelos de \ac{IA}/\ac{ML} y en el cumplimiento del marco normativo
aplicable:

\begin{enumerate}
  \item \textbf{Las características inherentes de \ac{ISO}/IEC~25012}
    ---completitud, exactitud, consistencia y actualidad--- implementadas a
    través de cuatro métricas (\texttt{completeness}, \texttt{syntactic\_accuracy},
    \texttt{logical\_consistency} y \texttt{currentness}) y de la detección de
    duplicados (\texttt{uniqueness}, sub-dimensión Con-ML-1 de consistencia
    según \ac{ISO}/IEC~5259-2). Esto proporciona cobertura directa de los
    requisitos de \ac{ISO}/IEC~5259-2~\cite{iso5259} e \ac{ISO}/IEC~25024~\cite{iso25024}.

  \item \textbf{Equilibrio de clases} (\texttt{class\_balance}), medido
    mediante la entropía de Shannon. El desequilibrio de clases es uno de los
    factores documentados de sesgo en sistemas de \ac{IA}~\cite{amazon2018,
    buolamwini2018} y su corrección tiene un impacto directo en la equidad y
    fiabilidad de los modelos de clasificación.

  \item \textbf{Identificabilidad de \ac{PII}} mediante enmascaramiento
    automático de columnas sensibles. Su inclusión en el \ac{MVP} responde a
    la obligación legal establecida por el \ac{RGPD}~\cite{rgpd2016} y por el
    Reglamento Europeo de \ac{IA}~\cite{aiact2024} de proteger los datos
    personales utilizados en el entrenamiento de modelos.
\end{enumerate}

Adicionalmente, los valores atípicos (\emph{outliers}) se han incorporado como
herramienta de diagnóstico en el módulo de \ac{EDA}, sin ser computados como
parte de la puntuación de calidad global. Esta decisión sigue explícitamente
la recomendación de \ac{ISO}/IEC~5259-2~\cite{iso5259}, que no considera los
\emph{outliers} una dimensión de calidad inherente, sino un indicador diagnóstico cuya
relevancia depende del dominio de aplicación.

El resultado es una delimitación deliberada del alcance del proyecto en la
frontera donde el dato pasa a depender del sistema que lo aloja: \dataqual
ofrece una auditoría algorítmica y estadística de la calidad intrínseca del
dato para proyectos de \ac{IA}, y deja fuera de su ámbito las dimensiones
cuya evaluación requiere acceso a la infraestructura o a los metadatos
organizativos de la entidad propietaria.


%% =========================================================================
\section{Estado del arte: herramientas existentes}
\label{sec:estado-arte}
%% =========================================================================

En el ecosistema actual existen diversas herramientas orientadas a la
evaluación y validación de la calidad de datos. Esta sección analiza las más
relevantes, examinando sus fortalezas y limitaciones para contextualizar la
contribución de \dataqual.

\subsection{Great Expectations}
\label{sec:great-expectations}

Great Expectations~\cite{greatexpectations} es una biblioteca de código abierto
escrita en Python para la definición y validación de «expectativas»
(\emph{expectations}) sobre conjuntos de datos. Las expectativas son
aserciones declarativas que describen las propiedades que se esperan de los
datos, como «esta columna no debe contener valores nulos» o «los valores de
esta columna deben estar comprendidos entre 0 y 100».

\paragraph{Fortalezas.}
Great Expectations cuenta con un ecosistema maduro y una comunidad activa.
Soporta múltiples \emph{backends} de ejecución (Pandas, Spark, \ac{SQL}),
lo que permite evaluar datasets de distintos tamaños y procedencias. Además,
genera automáticamente documentación visual de los resultados de validación
(\emph{Data Docs}).

\paragraph{Limitaciones.}
Se trata de una herramienta \emph{code-first} que requiere conocimientos de
Python para su configuración y uso. No dispone de una interfaz gráfica
integrada, lo que excluye a usuarios no técnicos. No ofrece Quality Gates
nativos, no implementa un sistema de rastreo de issues entre ejecuciones
sucesivas y presenta una curva de aprendizaje pronunciada para equipos sin
experiencia en programación.

\subsection{Apache Deequ}
\label{sec:apache-deequ}

Apache Deequ~\cite{schelter2018} es una biblioteca de verificación de calidad
de datos desarrollada por Amazon y construida sobre Apache Spark. Está diseñada
para operar a gran escala sobre datasets masivos en entornos distribuidos.

\paragraph{Fortalezas.}
Su principal ventaja es la escalabilidad masiva que hereda de Spark, lo que
permite procesar datasets de millones de registros de forma distribuida.
Ofrece integración nativa con los servicios de Amazon Web Services, verificación
declarativa de restricciones y detección de anomalías basada en series
temporales.

\paragraph{Limitaciones.}
Requiere un entorno JVM y un clúster de Spark para su ejecución, lo que supone
una barrera de entrada considerable. No dispone de interfaz web, resulta
desproporcionado para datasets de tamaño medio y no implementa métricas
específicas para \ac{ML} como el equilibrio de clases.

\subsection{DQOps}
\label{sec:dqops}

DQOps~\cite{dqops} es una herramienta de calidad de datos de código abierto
inspirada en la filosofía de SonarQube~\cite{sonarqube}. Proporciona un enfoque
sistemático a la calidad de datos con monitorización continua y alertas
automatizadas.

\paragraph{Fortalezas.}
DQOps comparte una filosofía similar a la de \dataqual en cuanto al
tratamiento de la calidad de datos como un proceso continuo. Soporta múltiples
fuentes de datos (BigQuery, Snowflake, PostgreSQL), ofrece una interfaz web
para la configuración y visualización de resultados, y permite definir alertas
automatizadas ante degradaciones de calidad.

\paragraph{Limitaciones.}
Está orientado fundamentalmente a bases de datos y no a ficheros de datos
(datasets en formatos \ac{CSV}, \ac{XLSX} o Parquet). Su configuración es
compleja, no ha sido diseñado específicamente para proyectos de \ac{IA}/\ac{ML}
y carece de métricas como el equilibrio de clases o la detección de outliers.

\subsection{Herramientas comerciales}
\label{sec:herramientas-comerciales}

Existen también soluciones comerciales orientadas a la preparación de datos,
como Trifacta, Alteryx y Talend. Estas herramientas ofrecen interfaces gráficas
pulidas, capacidades de transformación de datos y soporte empresarial.

No obstante, presentan limitaciones significativas para el caso de uso de
\dataqual: su coste de licenciamiento es elevado, generan dependencia del
proveedor (\emph{vendor lock-in}), están orientadas a la preparación y
transformación de datos (\emph{data wrangling}) más que a la evaluación de
calidad para \ac{IA}, y muchas de ellas operan como servicios en la nube
(\ac{SaaS}), lo que compromete la soberanía de los datos.

\subsection{Análisis comparativo}
\label{sec:analisis-comparativo}

El Cuadro~\ref{tab:comparativa} presenta una comparación sistemática de
\dataqual con las tres herramientas de código abierto analizadas. Se han
seleccionado quince características que cubren la funcionalidad, la
accesibilidad, las capacidades específicas para \ac{IA}/\ac{ML} y los aspectos
operativos.

\begin{table}[htbp]
  \centering
  \caption{Comparativa de \dataqual con herramientas de calidad de datos existentes}
  \label{tab:comparativa}
  \small
  \begin{tabularx}{\textwidth}{X c c c c}
    \toprule
    \textbf{Característica} & \textbf{\dataqual} & \textbf{Great Exp.} & \textbf{Deequ} & \textbf{DQOps} \\
    \midrule
    Interfaz web                    & Sí  & No  & No  & Sí  \\
    Requiere programación           & No  & Sí  & Sí  & Parcial \\
    Métricas predefinidas           & 7   & 300+ & 20+ & 150+ \\
    Detección automática de tipos   & Sí  & No  & No  & Parcial \\
    Quality Gates                   & Sí  & No  & No  & Sí  \\
    Fingerprinting de issues        & Sí  & No  & No  & No  \\
    Comparación entre ejecuciones   & Sí  & No  & Sí  & Sí  \\
    Equilibrio de clases            & Sí  & No  & No  & No  \\
    Detección de outliers           & Sí  & Parcial & Sí & Parcial \\
    Enmascaramiento de \ac{PII}     & Sí  & No  & No  & No  \\
    Versionado de datasets          & Sí  & No  & No  & No  \\
    Procesamiento asíncrono         & Sí  & No  & Sí  & Sí  \\
    Autoalojable                    & Sí  & Sí  & Sí  & Sí  \\
    Específico para IA/ML           & Sí  & No  & No  & No  \\
    Coste                           & Gratuito & Gratuito & Gratuito & Freemium \\
    \bottomrule
  \end{tabularx}
\end{table}

Como se puede observar, \dataqual ocupa un nicho diferenciado al combinar una
interfaz web completa con métricas específicas para \ac{IA}/\ac{ML},
fingerprinting de issues y autoalojamiento. Si bien Great Expectations supera
ampliamente a \dataqual en el número de métricas predefinidas, su naturaleza
\emph{code-first} limita su accesibilidad. Apache Deequ ofrece una
escalabilidad superior, pero a costa de una complejidad operativa que lo hace
inviable para equipos pequeños o datasets de tamaño moderado. DQOps comparte
la filosofía de calidad continua, pero su orientación a bases de datos y la
ausencia de métricas específicas para \ac{ML} lo sitúan en un segmento
diferente.


%% =========================================================================
\section{Diferenciación de \dataqual}
\label{sec:diferenciacion}
%% =========================================================================

A partir del análisis comparativo anterior, se identifican siete factores
clave que diferencian a \dataqual de las herramientas existentes:

\begin{enumerate}
  \item \textbf{Solución \emph{full-stack} con interfaz integrada.} A diferencia de
    herramientas \emph{code-only} como Great Expectations o Apache Deequ,
    \dataqual proporciona una interfaz web completa construida con
    React~\cite{react} y Next.js~\cite{nextjs} que permite a usuarios no
    técnicos configurar métricas, lanzar evaluaciones y consultar resultados
    sin escribir código.

  \item \textbf{Diseño centrado en IA/ML.} \dataqual incorpora métricas
    específicas para proyectos de aprendizaje automático que no se encuentran
    en las herramientas generalistas: equilibrio de clases medido mediante
    entropía de Shannon y evaluación de la actualidad temporal de los datos
    (\emph{currentness}).

  \item \textbf{Sistema de fingerprinting SHA-256.} Inspirado en
    SonarQube~\cite{sonarqube}, \dataqual genera un hash \ac{SHA}-256 para cada
    issue detectado, lo que permite rastrear su ciclo de vida a lo largo de
    múltiples ejecuciones: issues nuevos, recurrentes y corregidos.

  \item \textbf{Quality Gates como puntos de decisión.} Cada evaluación produce
    un veredicto global (\texttt{PASSED}, \texttt{WARNING} o \texttt{FAILED})
    basado en umbrales configurables, proporcionando una señal inequívoca
    sobre la aptitud del dataset para su uso.

  \item \textbf{Protección de \ac{PII} integrada.} El enmascaramiento de
    columnas sensibles es una funcionalidad nativa de la plataforma, no un
    complemento externo. Esto facilita el cumplimiento del \ac{RGPD} y
    otras normativas de protección de datos.

  \item \textbf{Arquitectura de plugins extensible.} El motor de métricas se
    basa en una clase abstracta \texttt{BaseMetric} y un registro global
    (\texttt{METRIC\_REGISTRY}) que permiten añadir nuevas métricas sin
    modificar el código existente, siguiendo el principio abierto-cerrado
    de los patrones de diseño~\cite{design_patterns}.

  \item \textbf{Autoalojamiento con soberanía de datos.} \dataqual se despliega
    mediante Docker Compose~\cite{docker} sin dependencias de servicios
    en la nube. Todos los datos permanecen en la infraestructura del usuario,
    lo que garantiza la soberanía de los datos y facilita el cumplimiento
    normativo.
\end{enumerate}


%% =========================================================================
\section{Justificación académica}
\label{sec:justificacion-academica}
%% =========================================================================

El desarrollo de \dataqual integra de manera transversal múltiples áreas del
plan de estudios del Grado en Ingeniería Informática, lo que le confiere un
carácter integrador propio de un \ac{TFG}. A continuación se detalla la
correspondencia entre las áreas curriculares y los componentes de la
plataforma.

\subsection{Ingeniería del Software}

El proyecto aplica principios de arquitectura de software a gran escala:
arquitectura de microservicios orquestados con Docker
Compose~\cite{docker}, patrón \ac{MVC} en el backend
Flask~\cite{flask}, y patrones de diseño clásicos~\cite{design_patterns}
como Abstract Factory (registro de métricas) y Observer (sistema de eventos de
evaluación). Se ha seguido control de versiones con Git y metodologías ágiles
de desarrollo.

\subsection{Sistemas distribuidos}

La plataforma implementa procesamiento asíncrono mediante colas de mensajes
con Celery~\cite{celery} y Redis~\cite{redis} como \emph{broker} de mensajes.
La orquestación de los ocho servicios Docker incluye \emph{healthchecks},
reinicio automático y comunicación inter-servicio mediante redes internas.
Flower proporciona monitorización en tiempo real de los \emph{workers}
asíncronos.

\subsection{Bases de datos}

El esquema relacional en PostgreSQL~\cite{postgresql} incluye el uso de
campos \texttt{JSONB} para configuraciones flexibles de métricas, migraciones
gestionadas con Flask-Migrate/Alembic y almacenamiento de objetos compatible
con S3 mediante MinIO~\cite{minio}. El \ac{ORM} SQLAlchemy~\cite{sqlalchemy}
proporciona la capa de abstracción sobre la base de datos relacional.

\subsection{Desarrollo web}

\dataqual es una aplicación \emph{full-stack} que combina un \emph{frontend} construido con
React~\cite{react}, Next.js~\cite{nextjs}, TypeScript~\cite{typescript} y
Material \ac{UI}~\cite{materialui}, con un backend \ac{REST}ful
implementado en Flask~\cite{flask}. La autenticación se gestiona mediante
tokens \ac{JWT}, y la comunicación entre frontend y backend se realiza a
través de una \ac{API} \ac{REST}.

\subsection{Análisis de datos y estadística}

El motor de métricas implementa técnicas estadísticas como el rango
intercuartílico (\ac{IQR}) y el Z-score para la detección de outliers, la
entropía de Shannon para la evaluación del equilibrio de clases, y el
\ac{EDA} interactivo con histogramas, diagramas de caja, matrices de
correlación y detección de anomalías.

\subsection{Seguridad informática}

La plataforma incorpora múltiples capas de seguridad: autenticación y
autorización mediante \ac{JWT}, derivación de contraseñas con \ac{PBKDF2},
limitación de tasa de peticiones (\emph{rate limiting}), validación de
entradas en el servidor, y enmascaramiento de \ac{PII} para proteger datos
sensibles conforme al \ac{RGPD}.

\subsection{DevOps y despliegue}

La containerización con Docker~\cite{docker} y la orquestación con Docker
Compose permiten desplegar la plataforma completa con un único comando.
La configuración basada en variables de entorno facilita la adaptación a
distintos entornos (desarrollo, pruebas, producción), y Flower proporciona
un panel de monitorización para los \emph{workers} de Celery.

\bigskip

En conjunto, el desarrollo de \dataqual demuestra la capacidad de integrar
conocimientos de múltiples disciplinas informáticas en un producto de software
funcional, cohesionado y alineado con estándares internacionales, lo que
constituye el objetivo fundamental de un \ac{TFG} en Ingeniería Informática.


% Local Variables:
%  coding: utf-8
%  mode: latex
%  mode: flyspell
%  ispell-local-dictionary: "castellano8"
% End:
